\chapter{Projection and Well-Formedness}\label{ch:projection}

A global type is written from nobody's point of view, but it is played by roles that
each see only their own messages. \emph{Projection} computes what each role must do;
\emph{well-formedness} asks whether the roles, acting only on what they see, can
together produce exactly the runs the global type allows
\citep{honda2008multiparty,denielou2012communicating}. A protocol that is not
well-formed cannot be implemented faithfully by any set of processes, so there is no
point reducing it to CTL. LeanFM checks it first, with \texttt{session\_check}.

\section{Local types}

Projecting PlaceOrder onto each role keeps that role's sends and receives and turns
each choice into a \emph{select} (the role decides) or an \emph{offer} (the role is
told).

\begin{lstlisting}[language={},basicstyle=\ttfamily\footnotesize,breaklines=true]
-- Client
send CreateOrder to Shop;
offer {
  receive Quote from Shop;
  send Pay to Shop;
  offer {
    receive Receipt from Shop;
  } or {
    receive PaymentDeclined from Shop;
  }
} or {
  receive SoldOut from Shop;
}

-- Shop
receive CreateOrder from Client;
select {
  send Quote to Client;
  receive Pay from Client;
  send Charge to Bank;
  offer {
    receive Charged from Bank;
    send Receipt to Client;
  } or {
    receive Declined from Bank;
    send PaymentDeclined to Client;
  }
} or {
  send SoldOut to Client;
}

-- Bank
offer {
  receive Charge from Shop;
  select {
    send Charged to Shop;
  } or {
    send Declined to Shop;
  }
} or {
}
\end{lstlisting}


The bank's projection shows what projection is for. The bank takes no part in the
sold-out branch, so its local type offers that branch as empty: from the bank's point
of view the session may end without a word. The shop's projection is the only one with
a \texttt{select}: it alone decides. The check reports nothing to fix.

\begin{lstlisting}[language={},basicstyle=\ttfamily\footnotesize,breaklines=true]
PlaceOrder: no issues; 3 roles, 9 messages, 2 top-level calls
\end{lstlisting}


\section{What makes a protocol unplayable}

LeanFM refuses two kinds of protocol, and each refusal names the problem.

\paragraph{No one decides.} If two different roles start the branches of a choice,
nothing stops both from acting at once. The two roles race, and the run that results
may be in neither branch.

\begin{lstlisting}[language={},basicstyle=\ttfamily\footnotesize,breaklines=true]
protocol Race (Client, Shop) {
  choice {
    Cancel from Client to Shop;
  } or {
    Ship from Shop to Client;
  }
}
\end{lstlisting}

\begin{lstlisting}[language={},basicstyle=\ttfamily\footnotesize,breaklines=true]
no single role decides a choice, so it is a race: Client starts Cancel; Shop starts Ship
\end{lstlisting}


\paragraph{Someone cannot tell.} If a role behaves differently in two branches but
receives nothing that tells it which branch was taken, it must guess. Here the bank
should accept a capture in one branch and nothing in the other, yet no message tells
the bank which one it is in.

\begin{lstlisting}[language={},basicstyle=\ttfamily\footnotesize,breaklines=true]
protocol Blind (Client, Shop, Bank) {
  h: Hold from Client to Bank;
  choice {
    Capture from Shop to Bank;
    Ship from Shop to Client;
  } or {
    SoldOut from Shop to Client;
  }
}
\end{lstlisting}

\begin{lstlisting}[language={},basicstyle=\ttfamily\footnotesize,breaklines=true]
Bank cannot tell which branch Shop chose
\end{lstlisting}


\section{Why well-formedness comes first}

These two conditions are what make the reduction of Part~\ref{part:ctl} meaningful. The
CTL formulas of Chapter~\ref{ch:faithful} describe the global behaviour: which messages
may follow which. A system of processes produces global behaviour by each process
following its own local type. If the protocol is well-formed, the processes together
produce exactly the protocol's runs, and checking the system against the formulas
checks each process against its role. If it is not, the formulas describe a behaviour
no system of processes produces, and a faithful implementation does not exist.

Well-formedness here is the \emph{knowledge of choice} condition of multiparty session
types: every role that behaves differently after a choice learns which branch was
taken before it has to act \citep{honda2008multiparty}. LeanFM checks it syntactically
on the global type. Horlings and Jongmans check related sanity conditions as CTL over
the protocol's state space \citep{horlings2021dcjlint}; the reduction in this book
goes the other way, from the protocol to the whole of its behaviour as CTL.
